\begin{abstract}
This paper analyzes the COBIT 2019 framework in corporate Information and Technology (I\&T) Governance, marking 30 years of evolution by ISACA. It examines the distinction between governance and management, core principles, 5 architectural domains (EDM, APO, BAI, DSS, MEA), and 7 governance components. Acting as a strategic umbrella, COBIT orchestrates frameworks such as ITIL, ISO/IEC 27001, TOGAF, and PMBOK. Supported by Design Factors and the CMMI capability model, it customizes adaptive governance systems. Finally, an empirical case study of the Central Bank of Nigeria (CBN) demonstrates how shared governance and integrated risk reporting transform operational firefighting into sustainable corporate value.
\end{abstract}

\begin{resumo}
Este artigo analisa o framework COBIT 2019 na Governança Corporativa de Informação e Tecnologia (I\&T), marco de 30 anos de evolução pela ISACA. Examina-se a distinção entre governança e gestão, princípios fundamentais, os 5 domínios arquiteturais (EDM, APO, BAI, DSS, MEA) e os 7 componentes habilitadores. Como guarda-chuva estratégico, o COBIT orquestra modelos como ITIL, ISO/IEC 27001, TOGAF e PMBOK. Por meio dos Fatores de Design e da escala de capacidade CMMI, demonstra-se a personalização de sistemas adaptativos. Por fim, o caso do Banco Central da Nigéria (CBN) ilustra como a governança compartilhada e a gestão de riscos transformam operações reativas em valor corporativo sustentável.
\end{resumo}
